% Code-Journal: je ein tatsächlicher Übergang zwischen zwei Shader-Ständen.
% Bearbeitbare Lehrmaterialien; Quellzeilen stammen aus den angegebenen Dateien.
% Farbig: ersetzte/neue Zeilen. Grau: unveränderter Kontext.

\newcommand{\schritteDemoEins}{%
% transition: 00-konstante-farbe/START.glsl -> 01-violett.glsl
\begin{journalschritt}{Weiß wird Violett}
\tagline{Demo 1 · Übergang 1/3}\hfill{\fontsize{9}{11}\selectfont\color{muted}Grau: unveränderter Kontext}\par\vspace{1mm}
\schrittcode{vorher}{Bisher · \texttt{START.glsl}}{%
\schrittzeile{accent}{3}{0.75}{float rot = 1.0;}
\schrittzeile{accent}{4}{0.75}{float gruen = 1.0;}
\schrittzeile{accent}{5}{0.75}{float blau = 1.0;}
}
\vspace{1mm}
\schrittcode{nachher}{Jetzt · \texttt{01-violett.glsl}}{%
\schrittzeile{sea}{3}{0.75}{float rot = 0.45;}
\schrittzeile{sea}{4}{0.75}{float gruen = 0.32;}
\schrittzeile{sea}{5}{0.75}{float blau = 0.78;}
}
\schrittkommentar{Drei RGB-Werte ändern; die Farbausgabe bleibt erhalten.}{Jeder Kanal liegt zwischen 0 und 1. Alle Pixel bekommen dieselbe Farbe: Die Position wird hier noch nicht verwendet.}
\end{journalschritt}

% transition: 00-konstante-farbe/01-violett.glsl -> 02-koralle.glsl
\begin{journalschritt}{Eine andere Farbmischung}
\tagline{Demo 1 · Übergang 2/3}\hfill{\fontsize{9}{11}\selectfont\color{muted}Grau: unveränderter Kontext}\par\vspace{1mm}
\schrittcode{vorher}{Bisher · \texttt{01-violett.glsl}}{%
\schrittzeile{accent}{3}{0.75}{float rot = 0.45;}
\schrittzeile{muted}{4}{0.75}{float gruen = 0.32;}
\schrittzeile{accent}{5}{0.75}{float blau = 0.78;}
}
\vspace{1mm}
\schrittcode{nachher}{Jetzt · \texttt{02-koralle.glsl}}{%
\schrittzeile{sea}{3}{0.75}{float rot = 0.95;}
\schrittzeile{muted}{4}{0.75}{float gruen = 0.32;}
\schrittzeile{sea}{5}{0.75}{float blau = 0.20;}
}
\schrittkommentar{Rot erhöhen und Blau senken. Grün bleibt bei 0,32.}{Die drei Kanäle lassen sich unabhängig verändern. Vor dem Ausführen die neue Farbe vorhersagen.}
\end{journalschritt}

% transition: 00-konstante-farbe/02-koralle.glsl -> ZIEL.glsl
\begin{journalschritt}{Die Zeit verändert die Farbe}
\tagline{Demo 1 · Übergang 3/3}\hfill{\fontsize{9}{11}\selectfont\color{muted}Grau: unveränderter Kontext}\par\vspace{1mm}
\schrittcode{vorher}{Bisher · \texttt{02-koralle.glsl}}{%
\schrittzeile{accent}{3}{0.75}{float rot = 0.95;}
\schrittzeile{accent}{4}{0.75}{float gruen = 0.32;}
\schrittzeile{accent}{5}{0.75}{float blau = 0.20;}
}
\vspace{1mm}
\schrittcode{nachher}{Jetzt · \texttt{ZIEL.glsl}}{%
\schrittzeile{sea}{3}{0.75}{float rot = 0.5 + 0.5 * sin(iTime);}
\schrittzeile{sea}{4}{0.75}{float gruen = 0.5 + 0.5 * sin(iTime + 2.0);}
\schrittzeile{sea}{5}{0.75}{float blau = 0.5 + 0.5 * sin(iTime + 4.0);}
}
\schrittkommentar{Die drei Konstanten durch zeitabhängige Ausdrücke ersetzen.}{\texttt{iTime}: Sekunden. Sinus liefert $[-1,1]$; halbieren und um 0,5 verschieben ergibt $[0,1]$. $+2$ und $+4$ versetzen die Schwingungen.}
\end{journalschritt}

}

\newcommand{\schritteDemoZwei}{%
% transition: 01-pixelposition-und-normalisierung/START.glsl -> 01-normalisierung.glsl
\begin{journalschritt}{Pixelposition wird Farbe}
\tagline{Demo 2 · Übergang 1/2}\hfill{\fontsize{9}{11}\selectfont\color{muted}Grau: unveränderter Kontext}\par\vspace{1mm}
\schrittcode{vorher}{Bisher · \texttt{START.glsl}}{%
\schrittzeile{accent}{3}{0.75}{vec3 color = vec3(1.0);}
}
\vspace{1mm}
\schrittcode{nachher}{Jetzt · \texttt{01-normalisierung.glsl}}{%
\schrittzeile{sea}{3}{0.75}{vec2 uv = fragCoord / iResolution.xy;}
\schrittzeile{sea}{4}{0.75}{vec3 color = vec3(uv.x, uv.y, 0.0);}
}
\schrittkommentar{Normierte Koordinaten ergänzen und die Farbe aus x und y bilden.}{\texttt{fragCoord}: Pixelposition; \texttt{iResolution.xy}: Breite und Höhe. Komponentenweise teilen: links/unten nahe 0, rechts/oben nahe 1.}
\end{journalschritt}

% transition: 01-pixelposition-und-normalisierung/01-normalisierung.glsl -> ZIEL.glsl
\begin{journalschritt}{Zwischen zwei Farben mischen}
\tagline{Demo 2 · Übergang 2/2}\hfill{\fontsize{9}{11}\selectfont\color{muted}Grau: unveränderter Kontext}\par\vspace{1mm}
\schrittcode{vorher}{Bisher · \texttt{01-normalisierung.glsl}}{%
\schrittzeile{accent}{4}{0.75}{vec3 color = vec3(uv.x, uv.y, 0.0);}
}
\vspace{1mm}
\schrittcode{nachher}{Jetzt · \texttt{ZIEL.glsl}}{%
\schrittzeile{sea}{4}{0.75}{vec3 links = vec3(0.08, 0.15, 0.45);}
\schrittzeile{sea}{5}{0.75}{vec3 rechts = vec3(1.0, 0.38, 0.18);}
\schrittzeile{sea}{6}{0.75}{vec3 color = mix(links, rechts, uv.x);}
}
\schrittkommentar{Zwei Referenzfarben ergänzen; mit der x-Position mischen.}{$\operatorname{mix}(a,b,t)=(1-t)a+tb$. Links gilt $t\approx0$, rechts $t\approx1$; in der Mitte werden beide Farben gleich gewichtet.}
\end{journalschritt}

}

\newcommand{\schritteDemoDrei}{%
% transition: 03-kreis-und-quadrierter-abstand/START.glsl -> 01-kreis.glsl
\begin{journalschritt}{Ein Abstandstest zeichnet einen Kreis}
\tagline{Demo 3 · Übergang 1/4}\hfill{\fontsize{9}{11}\selectfont\color{muted}Grau: unveränderter Kontext}\par\vspace{1mm}
\schrittcode{vorher}{Bisher · \texttt{START.glsl}}{%
\schrittzeile{accent}{3}{0.75}{vec3 color = vec3(1.0);}
}
\vspace{1mm}
\schrittcode{nachher}{Jetzt · \texttt{01-kreis.glsl}}{%
\schrittzeile{sea}{3}{0.75}{vec2 p = (fragCoord - 0.5 * iResolution.xy) / iResolution.y;}
\schrittzeile{sea}{4}{0.75}{vec2 a = p;}
\schrittzeile{sea}{5}{0.75}{float abstandQuadrat = dot(a, a);}
\schrittzeile{sea}{6}{0.75}{float radius = 0.3;}
\schrittzeile{sea}{7}{0.75}{vec3 color = vec3(0.04, 0.08, 0.16);}
\schrittzeile{sea}{8}{0.75}{if (abstandQuadrat < radius * radius) \{}
\schrittzeile{sea}{9}{1.50}{color = vec3(0.2, 0.85, 0.7);}
\schrittzeile{sea}{10}{0.75}{\}}
}
\schrittkommentar{Koordinaten, quadrierten Abstand, Radius und die Farbentscheidung ergänzen.}{Bildmitte abziehen, beide Achsen durch die Höhe teilen. \texttt{a=p}: Mittelpunkt im Ursprung. $\operatorname{dot}(a,a)=a_x^2+a_y^2$; innen gilt $d^2<r^2$.}
\end{journalschritt}

% transition: 03-kreis-und-quadrierter-abstand/01-kreis.glsl -> 02-ein-abstandsfeld.glsl
\begin{journalschritt}{Aus Abstand wird ein Einfluss}
\tagline{Demo 3 · Übergang 2/4}\hfill{\fontsize{9}{11}\selectfont\color{muted}Grau: unveränderter Kontext}\par\vspace{1mm}
\schrittcode{vorher}{Bisher · \texttt{01-kreis.glsl}}{%
\schrittzeile{accent}{6}{0.75}{float radius = 0.3;}
\schrittzeile{muted}{7}{0.75}{vec3 color = vec3(0.04, 0.08, 0.16);}
\schrittzeile{accent}{8}{0.75}{if (abstandQuadrat < radius * radius) \{}
}
\vspace{1mm}
\schrittcode{nachher}{Jetzt · \texttt{02-ein-abstandsfeld.glsl}}{%
\schrittzeile{sea}{6}{0.75}{float feld = 0.035 / (abstandQuadrat + 0.01);}
\schrittzeile{muted}{7}{0.75}{vec3 color = vec3(0.04, 0.08, 0.16);}
\schrittzeile{sea}{8}{0.75}{if (feld > 1.0) \{}
}
\schrittkommentar{Radius-Zeile und Bedingung ersetzen; die Farben bleiben gleich.}{Größerer Abstand: kleinerer Einfluss. $+0{,}01$ verhindert Division durch null; Maßstab und Schwelle bestimmen die Größe. Der Radius muss nicht erhalten bleiben.}
\end{journalschritt}

% transition: 03-kreis-und-quadrierter-abstand/02-ein-abstandsfeld.glsl -> 03-metaballs-statisch.glsl
\begin{journalschritt}{Ein zweites Feld kommt hinzu}
\tagline{Demo 3 · Übergang 3/4}\hfill{\fontsize{9}{11}\selectfont\color{muted}Grau: unveränderter Kontext}\par\vspace{1mm}
\schrittcode{vorher}{Bisher · \texttt{02-ein-abstandsfeld.glsl}}{%
\schrittzeile{accent}{4}{0.75}{vec2 a = p;}
\schrittzeile{muted}{5}{0.75}{float abstandQuadrat = dot(a, a);}
\schrittzeile{muted}{6}{0.75}{float feld = 0.035 / (abstandQuadrat + 0.01);}
}
\vspace{1mm}
\schrittcode{nachher}{Jetzt · \texttt{03-metaballs-statisch.glsl}}{%
\schrittzeile{sea}{4}{0.75}{float versatz = 0.23;}
\schrittzeile{sea}{5}{0.75}{vec2 a = p - vec2(-versatz, 0.0);}
\schrittzeile{sea}{6}{0.75}{vec2 b = p - vec2(versatz, 0.0);}
\schrittzeile{muted}{7}{0.75}{float abstandQuadrat = dot(a, a);}
\schrittzeile{muted}{8}{0.75}{float feld = 0.035 / (abstandQuadrat + 0.01);}
\schrittzeile{sea}{9}{0.75}{feld += 0.035 / (dot(b, b) + 0.01);}
}
\schrittkommentar{Zentren symmetrisch versetzen; den zweiten Feldbeitrag addieren.}{$p-M$ misst relativ zum Mittelpunkt. \texttt{+=} addiert zum bisherigen Wert. Die gemeinsame Schwelle wird erst auf die Summe angewendet: Formen können verschmelzen.}
\end{journalschritt}

% transition: 03-kreis-und-quadrierter-abstand/03-metaballs-statisch.glsl -> ZIEL.glsl
\begin{journalschritt}{Die Mittelpunkte bewegen sich}
\tagline{Demo 3 · Übergang 4/4}\hfill{\fontsize{9}{11}\selectfont\color{muted}Grau: unveränderter Kontext}\par\vspace{1mm}
\schrittcode{vorher}{Bisher · \texttt{03-metaballs-statisch.glsl}}{%
\schrittzeile{accent}{4}{0.75}{float versatz = 0.23;}
}
\vspace{1mm}
\schrittcode{nachher}{Jetzt · \texttt{ZIEL.glsl}}{%
\schrittzeile{sea}{4}{0.75}{float versatz = 0.23 + 0.14 * sin(iTime);}
}
\schrittkommentar{Nur den konstanten Versatz durch eine Schwingung ersetzen.}{0,23 ist die Mitte, 0,14 die Amplitude: Der Versatz liegt zwischen 0,09 und 0,37. Die Zentren bei $\pm\text{Versatz}$ bewegen sich gegenläufig.}
\end{journalschritt}

}

\newcommand{\schritteDemoVier}{%
% transition: 05-polynom-und-weiche-kante/START.glsl -> 01-zwei-punkte-lerp.glsl
\begin{journalschritt}{Eine weitere Stimme des Kanons}
\tagline{Demo 4 · Übergang 1/4}\hfill{\fontsize{9}{11}\selectfont\color{muted}Grau: unveränderter Kontext}\par\vspace{1mm}
\schrittcode{vorher}{Bisher · \texttt{START.glsl}}{%
\schrittzeile{accent}{23}{0.75}{for (int index = 0; index < 1; index += 2) \{}
}
\vspace{1mm}
\schrittcode{nachher}{Jetzt · \texttt{01-zwei-punkte-lerp.glsl}}{%
\schrittzeile{sea}{23}{0.75}{for (int index = 0; index < 4; index += 2) \{}
}
\vspace{1mm}
\schrittcode{basis}{Schon im START · Zeitversatz und gemeinsame Kurve erläutern}{%
\schrittzeile{muted}{24}{1.50}{float punktzeit = mod(iTime - float(index) * 1.5, 4.0 *}
\schrittzeile{muted}{}{2.25}{abschnittsdauer);}
\schrittzeile{muted}{48}{1.50}{vec2 punkt = mix(a, b, verlauf(t));}
\schrittzeile{muted}{52}{1.50}{vec3 punktfarbe = mix(vec3(0.0), eckfarbe, verlauf(puls));}
}
\schrittkommentar{Nur die Schleifengrenze von 1 auf 4 erhöhen.}{Schrittweite 2 wählt $j=0,2$: Stimmen 1 und 3, mit 0/3\,s Versatz. START-Rahmen erläutern: Zeitplan, Eckpunkte, Abstandstest und Farbimpuls im Halt.}
\end{journalschritt}

% transition: 05-polynom-und-weiche-kante/01-zwei-punkte-lerp.glsl -> 02-zwei-punkte-smoothstep.glsl
\begin{journalschritt}{Sanft bremsen und anfahren}
\tagline{Demo 4 · Übergang 2/4}\hfill{\fontsize{9}{11}\selectfont\color{muted}Grau: unveränderter Kontext}\par\vspace{1mm}
\schrittcode{vorher}{Bisher · \texttt{01-zwei-punkte-lerp.glsl}}{%
\schrittzeile{accent}{4}{0.75}{return t;}
}
\vspace{1mm}
\schrittcode{nachher}{Jetzt · \texttt{02-zwei-punkte-smoothstep.glsl}}{%
\schrittzeile{sea}{4}{0.75}{return t * t * (3.0 - 2.0 * t);}
}
\schrittkommentar{Nur die Rückgabe von verlauf durch das kubische Polynom ersetzen.}{$s(t)=3t^2-2t^3$: gleicher Anfang und gleiches Ende, dort flache Steigung. Dieselbe Kurve glättet Weg und Farbe. Die Haltezeit kommt weiterhin aus dem Zeitplan.}
\end{journalschritt}

% transition: 05-polynom-und-weiche-kante/02-zwei-punkte-smoothstep.glsl -> 03-vier-punkte-smoothstep.glsl
\begin{journalschritt}{Den Kanon vervollständigen}
\tagline{Demo 4 · Übergang 3/4}\hfill{\fontsize{9}{11}\selectfont\color{muted}Grau: unveränderter Kontext}\par\vspace{1mm}
\schrittcode{vorher}{Bisher · \texttt{02-zwei-punkte-smoothstep.glsl}}{%
\schrittzeile{accent}{23}{0.75}{for (int index = 0; index < 4; index += 2) \{}
}
\vspace{1mm}
\schrittcode{nachher}{Jetzt · \texttt{03-vier-punkte-smoothstep.glsl}}{%
\schrittzeile{sea}{23}{0.75}{for (int index = 0; index < 4; index += 1) \{}
}
\schrittkommentar{Nur die Schrittweite von 2 auf 1 ändern.}{Jetzt werden $j=0,1,2,3$ ausgeführt. Die zusätzlichen Stimmen erhalten 1,5/4,5\,s Versatz. Kurve, Bahn, Haltezeit und Farbwechsel bleiben gleich.}
\end{journalschritt}

% transition: 05-polynom-und-weiche-kante/03-vier-punkte-smoothstep.glsl -> ZIEL.glsl
\begin{journalschritt}{Den Verlauf weiter glätten}
\tagline{Demo 4 · Übergang 4/4}\hfill{\fontsize{9}{11}\selectfont\color{muted}Grau: unveränderter Kontext}\par\vspace{1mm}
\schrittcode{vorher}{Bisher · \texttt{03-vier-punkte-smoothstep.glsl}}{%
\schrittzeile{accent}{4}{0.75}{return t * t * (3.0 - 2.0 * t);}
}
\vspace{1mm}
\schrittcode{nachher}{Jetzt · \texttt{ZIEL.glsl}}{%
\schrittzeile{sea}{4}{0.75}{return t * t * t * (t * (6.0 * t - 15.0) + 10.0);}
}
\schrittkommentar{Nur die Rückgabe durch das vorbereitete Polynom fünften Grades ersetzen.}{Ausmultipliziert: $6t^5-15t^4+10t^3$. Die Beschleunigung endet ebenfalls bei null. Vertiefung: $s''(0)=s''(1)=0$; eine vollständige Herleitung ist optional.}
\end{journalschritt}

}

\newcommand{\schritteDemoFuenf}{%
% transition: 08-value-noise/START.glsl -> 01-zufall-pro-pixel.glsl
\begin{journalschritt}{Jedes Pixel erhält einen Zufallswert}
\tagline{Demo 5 · Übergang 1/3}\hfill{\fontsize{9}{11}\selectfont\color{muted}Grau: unveränderter Kontext}\par\vspace{1mm}
\schrittcode{vorher}{Bisher · \texttt{START.glsl}}{%
\schrittzeile{accent}{6}{0.75}{vec3 color = vec3(uv.x, uv.y, 0.0);}
}
\vspace{1mm}
\schrittcode{nachher}{Jetzt · \texttt{01-zufall-pro-pixel.glsl}}{%
\schrittzeile{sea}{8}{0.75}{float wert = zufall(floor(fragCoord));}
\schrittzeile{sea}{9}{0.75}{vec3 color = vec3(wert);}
}
\vspace{1mm}
\schrittcode{basis}{Vorbereitet ergänzen · Hash deklarieren und definieren}{%
\schrittzeile{ink}{1}{0.00}{float zufall(vec2 ort);}
\schrittzeile{ink}{16}{0.00}{float zufall(vec2 ort)}
\schrittzeile{ink}{17}{0.00}{\{}
\schrittzeile{ink}{18}{0.75}{return fract(sin(dot(ort, vec2(127.1, 311.7))) * 43758.5453);}
\schrittzeile{ink}{19}{0.00}{\}}
}
\schrittkommentar{Zufallswert abfragen und als drei gleiche Farbkanäle ausgeben.}{\texttt{floor} wählt die ganze Pixeladresse. Hash-Bausteine: Koordinaten gewichten, Sinus, vergrößern, Nachkommaanteil. Gleicher Ort liefert denselben Wert; ohne Zeit kein Flimmern.}
\end{journalschritt}

% transition: 08-value-noise/01-zufall-pro-pixel.glsl -> 02-value-noise.glsl
\begin{journalschritt}{Zwischen Gitterwerten interpolieren}
\tagline{Demo 5 · Übergang 2/3}\hfill{\fontsize{9}{11}\selectfont\color{muted}Grau: unveränderter Kontext}\par\vspace{1mm}
\schrittcode{vorher}{Bisher · \texttt{01-zufall-pro-pixel.glsl}}{%
\schrittzeile{accent}{8}{0.75}{float wert = zufall(floor(fragCoord));}
}
\vspace{1mm}
\schrittcode{nachher}{Jetzt · \texttt{02-value-noise.glsl}}{%
\schrittzeile{sea}{9}{0.75}{float wert = valueNoise(p);}
}
\vspace{1mm}
\schrittcode{basis}{Vorbereitet ergänzen und erklären · valueNoise, Funktionskörper}{%
\schrittzeile{ink}{25}{0.75}{vec2 zelle = floor(p);}
\schrittzeile{ink}{26}{0.75}{vec2 t = fract(p);}
\schrittzeile{ink}{28}{0.75}{vec2 weich = t * t * (3.0 - 2.0 * t);}
\schrittzeile{ink}{29}{0.75}{float unten = mix(zufall(zelle), zufall(zelle + vec2(1.0, 0.0)),}
\schrittzeile{ink}{}{1.50}{weich.x);}
\schrittzeile{ink}{30}{0.75}{float oben = mix(zufall(zelle + vec2(0.0, 1.0)), zufall(zelle +}
\schrittzeile{ink}{}{1.50}{vec2(1.0, 1.0)), weich.x);}
\schrittzeile{ink}{31}{0.75}{return mix(unten, oben, weich.y);}
}
\schrittkommentar{Aufruf wechseln; valueNoise deklarieren und ergänzen.}{\texttt{floor}: Zelle; \texttt{fract}: Zellposition. Vier Eckwerte horizontal, dann vertikal mischen; Smoothstep glättet.}
\end{journalschritt}

% transition: 08-value-noise/02-value-noise.glsl -> ZIEL.glsl
\begin{journalschritt}{Die Noise-Basis austauschen}
\tagline{Demo 5 · Übergang 3/3}\hfill{\fontsize{9}{11}\selectfont\color{muted}Grau: unveränderter Kontext}\par\vspace{1mm}
\schrittcode{vorher}{Bisher · \texttt{02-value-noise.glsl}}{%
\schrittzeile{accent}{9}{0.75}{float wert = valueNoise(p);}
}
\vspace{1mm}
\schrittcode{nachher}{Jetzt · \texttt{ZIEL.glsl}}{%
\schrittzeile{sea}{10}{0.75}{float wert = simplexNoise01(p);}
}
\vspace{1mm}
\schrittcode{basis}{Vorbereitet ergänzen · Simplex-Basis mit normierter Schnittstelle}{%
\schrittzeile{ink}{3}{0.00}{float simplexNoise01(vec2 p);}
\schrittzeile{ink}{111}{0.00}{float simplexNoise01(vec2 p)}
\schrittzeile{ink}{112}{0.00}{\{}
\schrittzeile{ink}{113}{0.75}{return clamp(0.5 + 0.5 * simplexNoise(p), 0.0, 1.0);}
\schrittzeile{ink}{114}{0.00}{\}}
}
\schrittkommentar{Im Gestaltungs-Code nur den Funktionsnamen austauschen.}{Simplex-Basis samt Lizenz übernehmen. Grundidee: Dreieck bestimmen, Richtungen auswerten, nach Abstand gewichten, Beiträge addieren. Beide Aufrufe liefern $[0,1]$; Maßstab und Farbe bleiben gleich.}
\end{journalschritt}

}

\newcommand{\schritteDemoSechs}{%
% transition: 09-fbm/START.glsl -> 01-fbm.glsl
\begin{journalschritt}{Grobe und feine Muster überlagern}
\tagline{Demo 6 · Übergang 1/3}\hfill{\fontsize{9}{11}\selectfont\color{muted}Grau: unveränderter Kontext}\par\vspace{1mm}
\schrittcode{vorher}{Bisher · \texttt{START.glsl}}{%
\schrittzeile{accent}{10}{0.75}{float wert = simplexNoise01(p);}
}
\vspace{1mm}
\schrittcode{nachher}{Jetzt · \texttt{01-fbm.glsl}}{%
\schrittzeile{sea}{6}{0.00}{float fbm(vec2 p)}
\schrittzeile{sea}{7}{0.00}{\{}
\schrittzeile{sea}{9}{0.75}{float wert = 0.5 * simplexNoise01(p)}
\schrittzeile{sea}{10}{2.81}{+ 0.25 * simplexNoise01(2.0 * p)}
\schrittzeile{sea}{11}{2.81}{+ 0.125 * simplexNoise01(4.0 * p)}
\schrittzeile{sea}{12}{2.81}{+ 0.0625 * simplexNoise01(8.0 * p);}
\schrittzeile{sea}{13}{0.75}{return wert / 0.9375;}
\schrittzeile{sea}{14}{0.00}{\}}
\schrittzeile{sea}{20}{0.75}{float wert = fbm(p);}
}
\schrittkommentar{Oktavsumme ergänzen und den Aufruf auf fbm umstellen.}{Basis aus Demo 5 unverändert. Größere Argumente erzeugen kleinere Strukturen. Gewichte halbieren; ihre Summe 0,9375 normalisiert den Mittelwert. Alle Oktaven liefern $[0,1]$.}
\end{journalschritt}

% transition: 09-fbm/01-fbm.glsl -> 02-wolken.glsl
\begin{journalschritt}{Aus Grauwerten werden Wolken}
\tagline{Demo 6 · Übergang 2/3}\hfill{\fontsize{9}{11}\selectfont\color{muted}Grau: unveränderter Kontext}\par\vspace{1mm}
\schrittcode{vorher}{Bisher · \texttt{01-fbm.glsl}}{%
\schrittzeile{accent}{21}{0.75}{vec3 color = vec3(wert);}
}
\vspace{1mm}
\schrittcode{nachher}{Jetzt · \texttt{02-wolken.glsl}}{%
\schrittzeile{sea}{21}{0.75}{float dichte = smoothstep(0.4, 0.65, wert);}
\schrittzeile{sea}{22}{0.75}{vec3 himmel = vec3(0.15, 0.4, 0.75);}
\schrittzeile{sea}{23}{0.75}{vec3 wolke = vec3(1.0);}
\schrittzeile{sea}{24}{0.75}{vec3 color = mix(himmel, wolke, dichte);}
}
\schrittkommentar{Wolkendichte und Farben ergänzen; die Graustufenausgabe ersetzen.}{Unter 0,4 bleibt der Himmel blau, über 0,65 wird er weiß. Dazwischen wirken Begrenzung, Normierung und das bekannte kubische Polynom; \texttt{mix} mischt die Farben.}
\end{journalschritt}

% transition: 09-fbm/02-wolken.glsl -> ZIEL.glsl
\begin{journalschritt}{Das Wolkenfeld wandert}
\tagline{Demo 6 · Übergang 3/3}\hfill{\fontsize{9}{11}\selectfont\color{muted}Grau: unveränderter Kontext}\par\vspace{1mm}
\schrittcode{vorher}{Bisher · \texttt{02-wolken.glsl}}{%
\schrittzeile{muted}{19}{0.75}{vec2 p = fragCoord / iResolution.y * 6.0;}
\schrittzeile{muted}{20}{0.75}{float wert = fbm(p);}
}
\vspace{1mm}
\schrittcode{nachher}{Jetzt · \texttt{ZIEL.glsl}}{%
\schrittzeile{muted}{19}{0.75}{vec2 p = fragCoord / iResolution.y * 6.0;}
\schrittzeile{sea}{20}{0.75}{p -= vec2(0.18, 0.06) * iTime;}
\schrittzeile{muted}{21}{0.75}{float wert = fbm(p);}
}
\schrittkommentar{Eine Zeile zwischen Koordinatenberechnung und Noise-Abfrage einfügen.}{Abfrageposition $p-vt$: Dasselbe Muster wandert nach rechts oben. Die Pixel bleiben an ihren Orten und berechnen neue Farben. Es wird kein neuer Zufall pro Bild erzeugt.}
\end{journalschritt}

}
